% JAC1964-20261121. Root deduction from H9-T7 and classical Hodge theory.
\subsection*{H9-T8a. Every fixed axis in characteristic zero}

\begin{corollary}[H9-T8a]\label{h9:T8a}
In H9-T7, suppose that the algebraically closed ground field $k$
has characteristic zero. There is no restriction on $m$ or on the
smooth axis of the genuine Lefschetz pencil. The signed height on
the entire group $\operatorname{CH}^m(Y)_{\rm hom,\mathbf Q}$ is
positive semidefinite, and its radical is exactly $\mathcal H_0$.
Thus the remaining arbitrary-axis sign question in H9-T7 is a
positive-characteristic question. This corollary uses classical
Hodge--Riemann positivity and does not require the new032 input.
\end{corollary}

\begin{proof}
First work over $\mathbf C$. The axis has dimension $2m-2$.
H9-T7 proves that every $w\in A_B$ is primitive for $h_B$.
It is the class of a rational algebraic cycle, so it is of Hodge
type $(m-1,m-1)$. The Hodge--Riemann bilinear relations give
\[
             (-1)^{m-1}\int_B w^2>0\qquad(w\ne0).
\]
On the ambient $2m$-fold the same relations give
$(-1)^m\int_X z^2>0$ for every nonzero primitive algebraic
middle class $z$. These are precisely the two forms in (H9.18).
Their orthogonal direct sum is positive definite, so H9-T7 gives
the asserted semipositivity and exact radical.

For a general algebraically closed $k$ of characteristic zero,
any proposed nonzero rational algebraic class and all the data
entering its intersection number descend to a finitely generated
field $k_1/\mathbf Q$. Choose an embedding $k_1\hookrightarrow
\mathbf C$. Geometric $\ell$-adic cohomology is unchanged by
algebraically closed extension, and the comparison with Betti
cohomology preserves cup products, algebraic cycle classes,
Lefschetz operators and traces. Thus the primitive condition and
the nonvanishing of the chosen class persist in the complex
realization, while its rational intersection number is unchanged.
The complex inequality therefore proves the assertion over $k$.
Applying this to every nonzero vector proves definiteness on both
algebraic subspaces, and the same isometry proves the radical
statement.
\end{proof}

\paragraph{Exact source and scope.}
The classical input is Griffiths, \emph{Periods of integrals on
algebraic manifolds, III}, Publ. Math. IHES 38 (1970), Section 2,
primitive space and intersection form (2.3)$'$ on printed p.~132,
conditionVI, (2.7), on printed p.~133;
\url{https://www.numdam.org/item/PMIHES_1970__38__125_0.pdf}.
Only its application to the two explicit summands in H9-T7 is
made here. This is not a claim of historical originality. In
Grothendieck's programme it supplies the fixed-axis form required
by cote44 scan4 after the admissible correction, using the
primitive convention of cote44 scan8 and cote16 no.7 scans34--40.
The proof concerns the original pencil field. A new total space
after a separable cover requires its own argument, as in H9-T12.
